\documentclass[12pt]{article}

\setlength{\parindent}{0pt}
\setcounter{secnumdepth}{3}
\usepackage[margin=1in]{geometry}

% packages
\usepackage{amsmath, amsfonts, amsthm, amssymb}
\usepackage{graphicx, float}
\usepackage{mathtools}
\usepackage{titlesec}
\usepackage{interval}
\usepackage{hyperref}
\usepackage{siunitx}
\usepackage{titling}
\usepackage{vwcol}
\usepackage{setspace}
\usepackage{empheq}
\usepackage{cancel}
\usepackage{esdiff}
\usepackage{multicol}
\usepackage{mdframed}
\usepackage{esdiff}
\usepackage{tikzsymbols}
\usepackage{multicol}
\usepackage{tikz}
\usepackage{varwidth}
\usepackage{pgfplots}
\usepackage{fancyhdr}
\usepackage{enumitem}
\usepackage{tcolorbox}

\intervalconfig {
	soft open fences
}

% header
\pagestyle{fancy}
\fancyhf{}
\lhead{Ethan Fan}
\rhead{AP Calculus BC}
\cfoot{\thepage}

% section format
\titleformat{\section}
  {\large \bfseries}
  {}
  {0em}
  {}
\titlespacing*{\section}{0pt}{0.5em}{0.3em}

\titleformat{\subsection}[runin]
  {\bfseries}
  {}
  {0em}
  {}
\titlespacing*{\subsection}{0pt}{0pt}{0em}

\titleformat{\subsubsection}[runin]
  {\itshape}
  {(\roman{subsubsection})}
  {0em}
  {}

% commands
\newcommand{\alignedintertext}[1]{%
  \noalign{%
    \vskip\belowdisplayshortskip
    \vtop{\hsize=\linewidth#1\par
    \expandafter}%
    \expandafter\prevdepth\the\prevdepth
  }%
}

\newcommand*{\paren}[1]{\ensuremath\left(#1\right)}
\newcommand*{\limit}[2][x]{\ensuremath{\displaystyle\lim_{#1\to#2}}}
\newcommand*{\Limit}[3][x]{\ensuremath{\displaystyle\lim_{#1\to#2}\left[#3\right]}}
\newcommand*{\deriv}[1][x]{\ensuremath{\dfrac{\mathrm{d}}{\mathrm{d}#1}}}
\newcommand*{\Deriv}[2][x]{\ensuremath{\dfrac{\mathrm{d}}{\mathrm{d}#1}\left[#2\right]}}
\newcommand{\dydx}{\dfrac{dy}{dx}}
\newcommand{\dydt}{\dfrac{dy}{dt}}
\newcommand{\dxdt}{\dfrac{dx}{dt}}
\newcommand*{\iinteg}[2][x]{\ensuremath{\displaystyle\int #2\;\mathrm{d}#1}}
\newcommand*{\dinteg}[4][x]{\ensuremath{\displaystyle\int_{#2}^{#3}#4\;\mathrm{d}#1}}
\newcommand*{\abs}[1]{\ensuremath{\left|#1\right|}}
\newcommand*{\eps}{\varepsilon}
\newcommand*{\real}{\ensuremath\mathbb{R}}
\newcommand*{\integer}{\ensuremath\mathbb{Z}}
\newcommand*{\posint}{\ensuremath\mathbb{Z^+}}
\newcommand*{\cbrt}[1]{\ensuremath{\sqrt[3]{#1}}}
\newcommand*{\lheqzero}{\ensuremath{\underset{\text{L'H}}{\overset{\left[\frac00\right]}{=}}}}
\newcommand*{\lheqinfty}{\ensuremath{\underset{\text{L'H}}{\overset{\left[\frac{\infty}{\infty}\right]}{=}}}}
\newcommand{\tor}{\text{ or }}
\newcommand{\floor}[1]{\lfloor #1 \rfloor}

\newcommand{\vem}{\vspace{0.5em}}
\newcommand{\example}[1]{\section{Example #1}}
\newcommand{\problem}[1]{\section{Problem #1}}
\newcommand{\subproblem}[1]{\subsection{(#1)} \,}

\DeclareMathOperator{\dne}{DNE}
\DeclareMathOperator{\sgn}{sgn}

\newcommand{\definition}[1]{\begin{tcolorbox}[colback=red!5!white,colframe=red!75!black,parbox=false] #1 \end{tcolorbox}}

\newcommand{\theorem}[2]{\begin{tcolorbox}[title={#1},colback=blue!5!white,colframe=blue!75!black,parbox=false] #2 \end{tcolorbox}}

\allowdisplaybreaks
% \postdisplaypenalty=100000

% document
\begin{document}

\vspace*{0cm}

% opening
\begin{center}
    {\Large \bfseries Problem Set 1} \\[0.5cm]
    {\large Distance, Displacement, and Anti-derivatives} \\[0.35cm]
    {\normalsize September 1, 2026}
\end{center}

% problems

\problem{5}

\subproblem{a} $D=\pi \cdot 1^2=\boxed{\pi \text{ ft}}$ \vem

\subproblem{b} $s(6)-s(2)= \displaystyle \int_2^6 v(t) \, dt=\boxed{2-\dfrac{\pi}{2} \text{ ft.}}$ The value measures the displacement.\vem

\subproblem{c} $2>\dfrac{\pi}{2} \implies \boxed{[4,6]}$\vem

\subproblem{d} $s(t)$ increases on $\boxed{(0,2) \cup (4,6)}$. Max: $\boxed{t=2,6}$. Min: $\boxed{t=4}$.

\subproblem{e} $\boxed{s(8)=0}$. Farthest: $\boxed{2 \text{ ft}}$\vem

\problem{6}

\subproblem{a} $A_A>A_B \implies$ rover A is ahead. \vem

\subproblem{b} When $v_A>v_B$, rover A is traveling at a greater speed and thus covering greater distances than rover B. The time at which the gap is largest is at $\boxed{t=4}$.\vem

\subproblem{c} 
\begin{align*}
    \int v_A(t) \, dt&= -\frac{1}{3}t^3+3t^2\\
    \int v_B(t) \, dt&= t^2
\end{align*}
\[
3(4)^2-\frac{1}{3}(4)^3-(4)^2=32-\frac{64}{3}=\boxed{\frac{32}{3} \text{ m}}
\]

\subproblem{d} 
\begin{gather*}
    t^2=-\frac{1}{3}t^3+3t^2\\
    \frac{1}{3}t^3=2t^2\\
    t^2(t-6)= 0\\
    \boxed{t=6}
\end{gather*}
Rover A came to a stop at $t=6$.\vem

\subproblem{e}
Both hold the same average velocity. $v_{A \, avg}=v_{B\, avg}=3$ m/s. \vem

\subproblem{f} (i): $t=4$. (ii): $t=2$. The largest gap takes place when the functions of the velocity intersect. The time when the gap is growing the fastest takes place when the functions of the velocity are parallel. 








\end{document}

